%% ==========================================================================
%%  Approach 14 --- almost-balanced size-shift route to the negative
%%  dichotomous subsidy problem.
%%
%%  Source: approach_14_proposal.pdf; full audit and closure in
%%  approach_14.md and RESIDUAL.md \S7.16.40--42.
%%
%%  STATUS.  CLOSED (negative). The size-shift reduction (Proposition 1 of
%%  the proposal) is correct, and its target (Research Question 1) turns out
%%  to be Approach~6's Target G-bal (\Cref{def:targetGbal}), restated in
%%  different language --- still open. The proposal's proof PROGRAM (modify
%%  the BKNS incremental algorithm to reach a good allocation) is refuted,
%%  by three independent obstructions, the last of which is a single
%%  hand-verified step requiring no search. Nothing here affects
%%  Conjecture~\ref{conj:main}, already proved unconditionally for every $n$
%%  by \Cref{sec:approach13}.
%% ==========================================================================
\section{Approach 14: Almost-Balanced Size-Shift (Closed, Negative)}
\label{sec:approach14}

\begin{quote}
\textbf{Closed, negative result.} This section records that the proof
\emph{program} proposed for Research Question~\ref{rq:g14-main} below cannot
succeed, and why. It does not close Research Question~\ref{rq:g14-main}
itself, which remains open and coincides with Approach~6's Target G-bal
(\Cref{def:targetGbal}). Full detail, all scripts, and every numeric witness
are in \texttt{approach\_14.md} at the repository root; this section gives
the mathematical content only.
\end{quote}

\subsection{The proposal}
\label{sec:g14-setup}

For negative dichotomous cost functions $\cost_i$ (marginals in $\set{0,1}$,
so $-\cost_i$ has marginals in $\set{0,-1}$), define the size-shifted
valuation $\tilde v_i(S) := \abs S - \cost_i(S)$. Every $\tilde v_i$ is a
positive dichotomous valuation. Writing
$w^{\cost}_{\alloc}(i,j) := \cost_i(\bundle i) - \cost_i(\bundle j)$ and
$\tilde w_{\alloc}(i,j) := \tilde v_i(\bundle j) - \tilde v_i(\bundle i)$, a
direct computation gives
\begin{equation}
\label{eq:g14-shift}
  w^{\cost}_{\alloc}(i,j) \;=\; \tilde w_{\alloc}(i,j) + \abs{\bundle i} - \abs{\bundle j}.
\end{equation}

Write $m = kn + r$ with $0 \le r < n$. An allocation $\alloc$ is
\emph{almost balanced} if $\abs{\bundle i} \in \set{k,k+1}$ for every $i$;
let $L$ be the agents with $\abs{\bundle i}=k$ and $H$ those with
$\abs{\bundle i}=k+1$.

\begin{definition}[Research Question 1]
\label{rq:g14-main}
For every positive dichotomous instance, is there an almost-balanced
allocation $\alloc$ and $\subsidy \in \set{0,1}^n$ such that $(\alloc,\subsidy)$
is envy-free and
\[
  \abs{\bundle i} > \abs{\bundle j} \implies \subsidy_i \le \subsidy_j
  \qquad \text{(i.e. } \subsidy_h \le \subsidy_\ell \text{ for all } h \in H,\ \ell \in L \text{)?}
\]
\end{definition}

\begin{proposition}
\label{prop:g14-reduction}
An affirmative answer to Question~\ref{rq:g14-main}, applied to
$\tilde v_i = \abs\cdot - \cost_i$, implies Conjecture~\ref{conj:main}.
\end{proposition}
\begin{proof}
Given such $(\alloc,\subsidy)$, set $d_i := \abs{\bundle i}-k \in \set{0,1}$
and $r_i := \subsidy_i + d_i \in \set{0,1,2}$. If $r_i=0$ then $i\in L$,
$\subsidy_i=0$; if $r_i=2$ then $i\in H$, $\subsidy_i=1$. Both values
occurring would give $\ell\in L$ with $\subsidy_\ell=0$ and $h\in H$ with
$\subsidy_h=1$, contradicting the hypothesis. So
$\max_i r_i - \min_i r_i \le 1$; with $\alpha:=\min_i r_i$ and
$p_i := r_i-\alpha \in \set{0,1}$, the shift-invariance of differences gives
$p_i-p_j = (\subsidy_i-\subsidy_j)+(\abs{\bundle i}-\abs{\bundle j})$, and by
\eqref{eq:g14-shift} and envy-freeness of $(\alloc,\subsidy)$,
$w^{\cost}_{\alloc}(i,j) \le p_i-p_j$ for all $i,j$. Thus $(\alloc,p)$ is an
envy-free solution of the original negative dichotomous instance with
$p\in\set{0,1}^n$.
\end{proof}

\begin{remark}
\label{rem:g14-notnew}
Question~\ref{rq:g14-main} is not a new statement: it is Approach~6's
\emph{Target G-bal} (\Cref{def:targetGbal}), which already requires a
cardinality-balanced partition (sizes differing by at most $1$) whose
optimal assignment gives $q$-spread at most $1$, where
$q_i = \tilde\subsidy_i + \abs{\bundle i}$. Since $q_i - k = r_i$ in the
notation above, the two statements coincide. Target G-bal was already known
to imply Conjecture~\ref{conj:main} (\Cref{obs:Gbal-implies-G} and
\Cref{prop:targetG-implies}) and remains open for $n \ge 3$; no
counterexample is known, including under the additional exhaustive search
performed for this section (\texttt{update\_14/rq1\_search.py}: zero
failures over almost-balanced allocations and all
$\subsidy\in\set{0,1}^n$, $n=3,4$).
\end{remark}

\subsection{The proposed proof program}
\label{sec:g14-program}

The proposal's route to Question~\ref{rq:g14-main} is to modify the BKNS
incremental algorithm (which builds an envy-freeable allocation with
$\subsidy\in\set{0,1}^n$ one item at a time via an \textsc{Extend} rule and a
\textsc{FindSink} rule) so that it additionally maintains, throughout
construction: (I1) envy-freeability with $\subsidy\in\set{0,1}^n$;
(I2) bundle sizes differing by at most $1$; (I3) every minimum-cardinality
agent lies in the maximally-subsidised set $M(\subsidy)$.

This program does not succeed, for three independent reasons.

\begin{proposition}[Whole-execution reachability]
\label{prop:g14-reach}
There is an explicit $3$-agent, $3$-item instance on which every legal
execution of the BKNS algorithm --- every item order, every valid
\textsc{Extend} choice, every \textsc{FindSink} starting point --- reaches
final bundle sizes $(0,1,2)$ and none reaches the almost-balanced profile
$(1,1,1)$, although some allocation on that instance achieves $q$-spread
$0$.
\end{proposition}
\begin{proof}[Verification]
Exhaustive backtracking over the full execution tree ($474$ nodes;
\texttt{update\_6/guidedR3\_full.py}, cross-checked by
\texttt{update\_14/reach\_sizes.py}) confirms the claim; an independent,
algorithm-free enumeration (\texttt{update\_6/verify\_reach\_gap.py})
confirms spread $0$ is achieved by the balanced allocation on the same
instance. Restricting \textsc{Extend} to minimum-cardinality bundles (Step 1
of the proposal) can only shrink the set of reachable allocations, and
(I2) is already unreachable on this instance.
\end{proof}

\begin{proposition}[Extend ignores cardinality]
\label{prop:g14-extend-ignores}
BKNS's \textsc{Extend} rule selects a receiving pair $(k,l)$ by
$l\in M(\subsidy)$ and a genuine marginal-$1$ gain
$v_k(g\mid \bundle l)=1$; the cardinality of $\bundle l$ plays no role.
Consequently there exist states in which the minimum-cardinality bundle lies
in $M(\subsidy)$ but \emph{every} valid \textsc{Extend} option grows a
strictly larger bundle.
\end{proposition}
\begin{proof}[Verification]
A minimal instance is exhibited in \texttt{update\_14/extend\_witness.py}:
$n=3$, partial allocation $\bundle{} = (\emptyset,\set 0,\set 4)$,
$\subsidy=(0,0,0)$, so $M(\subsidy)=\set{0,1,2}$; agent $0$ holds the unique
minimum-cardinality (empty) bundle and lies in $M(\subsidy)$, but
$v_0(3\mid\emptyset)=0$, so the only valid \textsc{Extend} option in the
entire state grows agent $1$'s strictly larger bundle. Frequency over random
instances is checked in \texttt{extend\_forced.py} and
\texttt{extend\_cardinality.py}: states where every valid option is off the
minimum occur regularly ($9$--$57$ times per few hundred trials,
$n=3,\dots,5$).
\end{proof}

\begin{theorem}[Closure: a balanced partition can still be hopeless]
\label{thm:g14-closure}
There is an explicit state, reachable by \textsc{Extend}'s own selection
rule, at which the resulting bundles are exactly balanced in cardinality yet
\emph{no} assignment of those bundles to the three agents achieves subsidy
spread at most $1$.
\end{theorem}
\begin{proof}
Take $n=3$ agents and bundles $\bundle 1=\set{b_1,b_2}$,
$\bundle2=\set{a_1,a_2}$, $\bundle3=\set{a_3}$ with
\[
  \tilde v_1(\bundle1,\bundle2,\bundle3) = (1,2,1), \quad
  \tilde v_2(\bundle1,\bundle2,\bundle3) = (0,2,1), \quad
  \tilde v_3(\bundle1,\bundle2,\bundle3) = (0,2,1).
\]
The identity assignment has welfare $4$, tied for maximal, hence
envy-freeable; its Halpern--Shah subsidy is $\subsidy=(1,0,1)$, so
$M(\subsidy)=\set{1,3}$, and agent $3$ holds the unique minimum-cardinality
bundle while lying in $M(\subsidy)$ --- exactly the configuration (I3) is
designed to produce.

Introduce a new item $c_1$ with $\tilde v_1(c_1\mid\bundle3)=1$ and
$\tilde v_2(c_1\mid\bundle3)=\tilde v_3(c_1\mid\bundle3)=0$. By
Proposition~\ref{prop:g14-extend-ignores}'s mechanism, the only valid
\textsc{Extend} witness for growing $\bundle3$ is agent $1$, not agent $3$.
After insertion the three bundles $\bundle1,\bundle2,\bundle3\cup\set{c_1}$
all have size $2$: cardinality balance (I2) is achieved. Yet direct
computation of all six assignments of these bundles to the three agents
gives subsidy spread exactly $2$ in every case, including the assignment
that leaves agent $3$ in place on its own enlarged bundle
($\subsidy=(2,0,1)$). No assignment of this triple of bundles to these three
agents achieves spread $\le 1$.
\end{proof}

\begin{corollary}
\label{cor:g14-closed}
The proof program of \Cref{sec:g14-program} --- restricting
\textsc{Extend}/\textsc{FindSink} choices to maintain (I1)--(I3) --- cannot
succeed. Weakening (I3) to the path condition
$\tilde d_{\alloc}(h,\ell)\le 0$ for $h\in H,\ell\in L$ (the proposal's
Step~4) does not rescue it: in the witness of
Theorem~\ref{thm:g14-closure}, (I3) is never violated, and the obstruction is
that spread $2$ is unavoidable on the given triple of bundles independent of
any subsidy-compatibility condition.
\end{corollary}
\begin{proof}
Immediate from Theorem~\ref{thm:g14-closure}: the state it exhibits is
reachable by \textsc{Extend}'s own rule (so no restriction of that rule can
avoid it without also making the rule inapplicable), and the failure is not
a choice among several outcomes but the total absence of a good outcome
among all six assignments of the resulting bundles.
\end{proof}

\subsection{Status}
\label{sec:g14-status}

Proposition~\ref{prop:g14-reduction} is correct and reusable: any future
route to Question~\ref{rq:g14-main} still yields Conjecture~\ref{conj:main}
via the same size-shift argument. Question~\ref{rq:g14-main} itself remains
open, coincides with Target G-bal, and is not attempted here beyond the
search of Remark~\ref{rem:g14-notnew}. What is closed is the specific
incremental-algorithm proof program: Proposition~\ref{prop:g14-reach},
Proposition~\ref{prop:g14-extend-ignores}, and Theorem~\ref{thm:g14-closure}
give three independent, mutually reinforcing obstructions, the last
requiring no search to verify. Any future attempt at
Question~\ref{rq:g14-main} must be able to move an item out of a bundle it
already occupies; insertion-only templates are excluded by this section.
None of this affects Conjecture~\ref{conj:main}, already proved
unconditionally for every $n$ in \Cref{sec:approach13}.
